\begin{deluxetable}{ccccc}
% \tablenum{1}
\tabletypesize{\scriptsize}
\tablewidth{0pt}
\tablecaption{Magnetic Field Strength \& Inclination Information \label{tab:magnetic_field_information}}
\tablehead{
\colhead{Dataset} & \colhead{Mean} & \colhead{Median} & \colhead{RMS} & \colhead{(Min, Max)}}
\startdata
DS0 & 8.5 G & 4.2 G & 24.6 G & (-738 G, 687 G) \\
{} & \nodata & \nodata & \nodata & \nodata \\
\hline
DS1 & 7.3 G & 5.3 G & 12.2 G & (-391 G, 287 G) \\
{} & \nodata & \nodata & \nodata & \nodata \\
\hline
DS2 & 102.7 G & 19.8 G & 238.7 G & (-1862 G, 880 G) \\
{} & 92.7\degree & 91.2\degree & 94.3\degree & (24\degree, 169\degree) \\
\hline
DS3 & 90.0 G & 12.7 G & 219.3 G & (-1694 G, 1165 G)\\
{} & 91.3\degree & 91.4\degree & 93.1\degree & (14\degree, 152\degree) \\
\hline
DS4 & 51.4 G & 9.4 G & 120.5 G & (-1182 G, 1042 G) \\
{} & 101.1\degree & 93.4\degree & 103.0\degree & (4\degree, 179\degree)
\enddata
\tablecomments{DS0 refers to the QS disk center ($\mu$ = 0.99) dataset taken on 25 April 2019 described in \citetalias{2023_Vesa}.}
\end{deluxetable}
